% !TEX root = ../main.tex

\phantomsection
\chapter*{초록}

\label{ch:abstract}

디파이(DeFi, 탈중앙 금융) 서비스(프로토콜)는 거래 상대를 가려내고 위험 매개변수(위험에 따라 정하는 설정값)를 정하려고 온체인(블록체인 위의) 지갑 평판을 찾아요. 적응형 가중 페이지랭크(AWP)는 토큰 송금(transfer)을 바탕으로 지갑에 순위를 매겨요. 하지만 송금은 돈을 낸 일을 기록할 뿐, 믿음을 기록하지는 않아요. ERC-20(토큰 공통 규칙) 허락(allowance, 토큰을 얼마까지 대신 써도 된다는 허락)은 주인(토큰 주인)이 스펜더(spender, 허락을 받아 대신 쓰는 쪽)에게 자기 토큰을 옮겨도 된다고 허락한 일을 기록해요. 이 논문은 페이지랭크(PageRank) 걷기가 송금과 함께 허락 화살표(엣지)도 따라가면 무엇을 얻는지 물어요.

이 논문의 방법은 결합 페이지랭크(C-PR, 두 그래프를 함께 걷는 점수)예요. C-PR은 한 번의 걷기 안에서, 모든 주소의 허락 줄과 송금 줄(그 주소에서 다음에 어디로 갈지 정하는 확률들)을 정해 둔 무게로 섞어요. C-PR의 한 층 비교 기준은 AWP와 엔도스랭크(EndorseRank)예요. 엔도스랭크는 0이 아닌 마지막 허락(가장 최근 허락) 위에서 AWP의 화살표 무게(가중치)를 쓰는 페이지랭크예요. 사용자 지갑은 너무 여러 가지로 쓰여서 한 잣대로 잴 수 없어요. 그래서 이 연구는 스마트 계약(블록체인 위에서 저절로 실행되는 프로그램)에 점수를 매겨요. 이 연구는 원시 차수(화살표를 그냥 센 수)를 포함한 모든 점수를, 동결일(점수를 고정하는 날)마다 그 앞 열두 달의 기록으로 계산해요. 대상은 계약 $6{,}483$개예요. 이 계약들은 2025년 7월부터 2026년 6월 사이에 아비트럼 원(Arbitrum One)에서, 다섯 토큰 가운데 하나로 적어도 세 주인에게서 허락을 받았어요. 이 연구는 점수들을 켄달의 타우($\tau$, 두 줄 세우기가 얼마나 같은 순서인지 나타내는 수)와 짝지은 부트스트랩(다시 뽑기) 구간으로 견주어요. 또 2026년 3월 31일에 동결한 점수를 스펜더 계약 $7{,}914$개에서 시험해요. 그리고 2026년 7월부터 9월까지의 기간을, 그 로그(계약이 남기는 기록)를 뽑기 전에 등록했어요. 등록 분석(미리 등록한 그대로 한 분석)의 결과를 안 뒤에, 이 연구는 금액을 달러(USD)로 다시 매기고, 다섯 토큰만 남기고, 열두 달 점수 기간을 정했어요. 이 연구의 주된 분석은 미리 정한 규칙을 이 고친 점수에 다시 적용해요. 그리고 결합 걷기에 대한 등록 분석의 판정도 함께 보여 줘요.

이 규칙들 아래에서 C-PR은 어느 한 층만 걷는 것도 이기지 못해요. C-PR은 나중의 허락 쌍(나중에 새로 생긴 주인–스펜더 허락)에서는 허락 층보다 못하지 않지만($\Delta\tau=+0.031$ $[0.021,0.042]$), 나중의 송금자(나중에 새로 송금을 보내온 주소)에서는 송금 층보다 낮아요($-0.030$ $[-0.040,-0.020]$). 그리고 등록 기간에서는 허락 쌍에서 이득을 보지만($+0.170$ $[0.149,0.189]$), 송금자에서는 허용 폭(더 나빠도 괜찮다고 미리 정한 폭)보다 더 많이 뒤처져요($-0.025$ $[-0.038,-0.015]$). 나중에(라벨을 본 뒤에) 계산한 대비(두 값의 차이)에서, C-PR은 두 기간 모두에서 두 라벨(나중에 채점에 쓰는 정답 값)을 엔도스랭크보다 더 잘 줄 세워요(허락 쌍에서 각각 $+0.026$ $[0.016,0.037]$, $+0.017$ $[0.006,0.029]$). 또 나중의 허락 쌍을 AWP보다 더 잘 줄 세워요($+0.133$ $[0.114,0.153]$). 이것은 주로 AWP가 0점을 주는 스펜더들 덕분이에요. 그리고 나중의 송금자는 AWP보다 못하게 줄 세워요($-0.073$ $[-0.083,-0.063]$). 이 대비들에서 어떤 페이지랭크도 나중에 새로 오는 것을 같은 종류 화살표를 그냥 센 수보다 더 잘 줄 세우지 못해요. 허락한 주인 수(서로 다른 주인 수)는 두 기간 모두에서 엔도스랭크를 이겨요(각각 $-0.090$ $[-0.100,-0.079]$, $-0.082$ $[-0.090,-0.073]$). AWP의 활동 재시작(활발한 주소에서 더 자주 다시 출발하기)을 쓴 엔도스랭크도 이겨요(각각 $-0.060$ $[-0.070,-0.052]$, $-0.046$ $[-0.055,-0.038]$). 그리고 송금 입차수(송금을 보내온 서로 다른 주소 수)는 AWP를 이겨요(각각 $-0.038$ $[-0.043,-0.032]$, $-0.041$ $[-0.047,-0.034]$). GMX~V2 계약 트레이더(GMX V2에서 거래하는 계약)가 나중에 청산(손실이 너무 커져 거래가 강제로 끝나는 일)을 피하는 일은 두 종류의 화살표 어느 쪽으로도 만들지 않은 결과 값이에요. 이 결과 값에서 허락한 주인 수는 등록 기간에 송금 수보다 앞섰어요($\Delta\tau=+0.141$, 97.5\% 구간(여러 번 비교할 때 넓힌 신뢰구간) $[0.026,0.256]$, 트레이더 $70$개). 하지만 더 이른 기간에서는 앞섰다고 가려지지 않았어요($+0.095$, 97.5\% 구간 $[-0.041,0.263]$). 이와 달리 등록 분석에서는 더 이른 기간에서만 가려졌어요. C-PR의 푸는 시간은 AWP의 $1.36$배예요. 엔도스랭크는 화살표가 $2.21$배 적은 그래프에서 돌고, 푸는 시간은 AWP의 $3.21$분의 1이에요.

이 논문은 다음을 기여(이 연구가 새로 보탠 것)로 내놓아요. 권한을 맡기는 일(허락)과 돈을 내는 일(송금)을 함께 묶는 걷기, 그리고 미리 정한 규칙 아래에서 그 걷기가 실패한 시험이에요. 평판 화살표로서의 허락이에요. 모든 점수를 그 점수가 매끄럽게 다듬는(평활하는) 수와 견준 평가예요. 이 평가는 같은 코호트(묶음)에서도, 기간 밖에서도 했어요. 그리고 한 걸음 깊이에서 이런 걷기가 따르는 닫힌 식(바로 계산할 수 있는 식)이에요. 이 식은 엔도스랭크를 주인 정규화 입강도(주인마다 나누어 맞춘, 들어오는 무게의 합)로 줄여요. 허락 그래프에서는 짜고 움직이는 주소 약 25개로 한 계약을 코호트의 90번째 백분위(100등분했을 때 아래에서 90번째 자리) 위로 올릴 수 있어요. 그리고 같은 한도 계산이 C-PR에도 들어맞아요. 그래서 시빌 공격(가짜 주소를 많이 만들어 속이는 공격)을 막는 장치 없이 어떤 점수로도 보상을 정해서는 안 돼요. 그리고 어떤 점수도 신용도(빌린 돈을 갚을 수 있는 정도)나 계약 코드의 안전을 재지 않아요.

\emph{핵심어:} 결합 페이지랭크; 탈중앙 금융(디파이); 온체인 평판; ERC-20 허락; 엔도스랭크; 적응형 가중 페이지랭크; 스마트 계약; 시간 홀드아웃; 미리 등록하기; 아비트럼 원(Arbitrum One); 켄달의 타우

% UNISA requires the title, the abstract and the key terms in two further official
% languages of South Africa. The translations are to be supplied by Unisa language
% services and inserted here after this page; they are not machine translated.
